\section{The Cylinder is a Metric Substitution, not an Inherited Geometry}
\label{sec:app_metric_choice}

\begin{lemma}[Amplitude--phase chart and its pullback metric]
\label{lem:chart_pullback}
Let $\mathbb{C}^* = \mathbb{C} \setminus \{0\}$ and let
\begin{equation}
    \Phi : \mathbb{C}^* \to (0, \infty) \times S^1, \qquad
    \Phi(z) = \left( |z|, \; z / |z| \right),
\end{equation}
with inverse $\Phi^{-1}(A, \theta) = (A\cos\theta, A\sin\theta)$. Then $\Phi$ is a diffeomorphism, and the pullback of the Euclidean metric $g_{\mathrm{euc}} = dx^2 + dy^2$ under $\Phi^{-1}$ is
\begin{equation}
    (\Phi^{-1})^* g_{\mathrm{euc}} = dA^2 + A^2 d\theta^2 .
\end{equation}
\end{lemma}
\begin{proof}
Both $\Phi$ and $\Phi^{-1}$ are smooth on their domains and are mutually inverse by construction, so $\Phi$ is a diffeomorphism. Differentiating $\Phi^{-1}$,
\begin{align}
    dx &= \cos\theta \, dA - A \sin\theta \, d\theta, \\
    dy &= \sin\theta \, dA + A \cos\theta \, d\theta .
\end{align}
Squaring and summing, the cross terms $\mp 2A\sin\theta\cos\theta \, dA \, d\theta$ cancel and the remaining terms group by $\cos^2\theta + \sin^2\theta = 1$:
\begin{equation}
    dx^2 + dy^2 = dA^2 + A^2 d\theta^2 . \qedhere
\end{equation}
\end{proof}

CyFM does not use this metric. It replaces it with the decoupled product metric $g_{\mathrm{cyl}} = dA^2 + d\theta^2$, that is, it deletes the factor $A^2$. The map $\Phi$ is therefore a diffeomorphism but \emph{not} an isometry between $(\mathbb{C}^*, g_{\mathrm{euc}})$ and $((0,\infty) \times S^1, g_{\mathrm{cyl}})$. This substitution, and not the identification of the domain, is the modelling decision this paper evaluates.

\begin{remark}[Flatness is not the distinguishing property]
\label{rem:both_flat}
Both metrics have identically vanishing Riemann curvature: $g_{\mathrm{euc}}$ because $\mathbb{C}^*$ is an open subset of the plane, and $g_{\mathrm{cyl}}$ because it is a product of one-dimensional factors. Each is therefore locally isometric to $\mathbb{R}^2$, and no local computation separates them. They are not globally isometric: under $g_{\mathrm{cyl}}$ the curves $A = \mathrm{const}$ are closed geodesics, whereas the geodesics of $g_{\mathrm{euc}}$ are straight line segments and none is closed. The non-trivial fundamental group $\pi_1 = \mathbb{Z}$ is a property of the punctured domain, shared by both, and is a cost the model pays through the $2\pi$ wrap rather than a benefit it gains.
\end{remark}

\begin{corollary}[Circumference]
\label{cor:circumference}
The circle $\{A = r\}$ has length $2\pi r$ under $g_{\mathrm{euc}}$ and length $2\pi$ under $g_{\mathrm{cyl}}$, for every $r > 0$.
\end{corollary}

This single difference accounts for both of the paper's velocity statements. Under $g_{\mathrm{euc}}$ the circumference available for a phase change shrinks to zero as $r \to 0$, so a path carrying a fixed phase difference through a neighbourhood of the origin must traverse it at ever greater angular speed; this is the $|c| / |z_t|^2$ growth of Proposition~\ref{thm:angular_divergence}. Under $g_{\mathrm{cyl}}$ the circumference is constant, so a half-turn costs the same at every amplitude and an angular displacement of $\pi$ always suffices; this is the bound of Lemma~\ref{lem:bounded_velocity}. The two results are the same fact about the factor $A^2$, read at the two ends of its range.

\section{Path-Level Consequences of the Euclidean Metric}
\label{sec:app_proofs}

We present the formal proofs for the statements made in the main text about what the inherited Euclidean metric on $(\mathrm{Re}, \mathrm{Im}) \in \mathbb{R}^2$ implies for probability paths and for the regression target.

Let $z_1 = e^{i\theta_1}$ and $z_2 = e^{i\theta_2}$ be two unit-amplitude phasors on the circle $S^1$, and let $\Delta\theta = \theta_1 - \theta_2 \in [-\pi, \pi]$ denote their phase difference.

\begin{lemma}[Chordal vs. Geodesic Distance]
\label{lem:chord_vs_arc}
The Euclidean (chordal) distance $d_{\mathrm{chord}}$ in the complex plane and the geodesic (arc) distance $d_{S^1}$ on the circle are related by:
\begin{align}
    d_{\mathrm{chord}}(\theta_1, \theta_2) &= \|e^{i\theta_1} - e^{i\theta_2}\|_2 = 2 \left|\sin\left(\frac{\Delta\theta}{2}\right)\right|, \\
    d_{S^1}(\theta_1, \theta_2) &= |\Delta\theta|.
\end{align}
\end{lemma}
\begin{proof}
Expanding the squared Euclidean distance:
\begin{align}
    \|e^{i\theta_1} - e^{i\theta_2}\|_2^2 &= (\cos\theta_1 - \cos\theta_2)^2 + (\sin\theta_1 - \sin\theta_2)^2 \\
    &= (\cos^2\theta_1 + \sin^2\theta_1) + (\cos^2\theta_2 + \sin^2\theta_2) - 2(\cos\theta_1\cos\theta_2 + \sin\theta_1\sin\theta_2) \\
    &= 2 - 2\cos(\theta_1 - \theta_2) = 2(1 - \cos\Delta\theta).
\end{align}
Using the half-angle trigonometric identity $1 - \cos\alpha = 2\sin^2(\alpha/2)$:
\begin{equation}
    \|e^{i\theta_1} - e^{i\theta_2}\|_2^2 = 4\sin^2\left(\frac{\Delta\theta}{2}\right).
\end{equation}
Taking the square root yields $d_{\mathrm{chord}} = 2 \left|\sin\left(\frac{\Delta\theta}{2}\right)\right|$. The geodesic distance on the unit circle is trivially the absolute arc length $|\Delta\theta|$.
\end{proof}


\begin{proposition}[Expected Midpoint Attenuation]
\label{prop:attenuation_proof}
Consider a linear (Euclidean) interpolation between two unit phasors: $z_t^{\mathrm{euc}} = (1-t)e^{i\theta_1} + t e^{i\theta_2}$ for $t \in [0, 1]$. Under a uniform phase difference prior $\Delta\theta \sim \mathcal{U}((-\pi, \pi])$, the expected magnitude at the midpoint $t=0.5$ is exactly $2/\pi$, representing an expected signal attenuation of:
\begin{equation}
    1 - \frac{2}{\pi} \approx 36.338\%.
\end{equation}
\end{proposition}
\begin{proof}
At $t=0.5$, the interpolated signal is:
\begin{equation}
    z_{0.5}^{\mathrm{euc}} = \frac{1}{2} (e^{i\theta_1} + e^{i\theta_2}) = e^{i\frac{\theta_1+\theta_2}{2}} \cos\left(\frac{\Delta\theta}{2}\right).
\end{equation}
The magnitude is $|z_{0.5}^{\mathrm{euc}}| = \left|\cos\left(\frac{\Delta\theta}{2}\right)\right|$. The expected value under a uniform distribution $p(\Delta\theta) = \frac{1}{2\pi}$ is:
\begin{equation}
    \mathbb{E}\left[|z_{0.5}^{\mathrm{euc}}|\right] = \frac{1}{2\pi} \int_{-\pi}^\pi \left|\cos\left(\frac{\Delta\theta}{2}\right)\right| d\Delta\theta.
\end{equation}
Since $\Delta\theta \in [-\pi, \pi]$, the half-angle $\Delta\theta/2 \in [-\pi/2, \pi/2]$. In this domain, the cosine function is non-negative, so the absolute value can be dropped:
\begin{align}
    \mathbb{E}\left[|z_{0.5}^{\mathrm{euc}}|\right] &= \frac{1}{2\pi} \int_{-\pi}^\pi \cos\left(\frac{\Delta\theta}{2}\right) d\Delta\theta \\
    &= \frac{1}{2\pi} \left[ 2\sin\left(\frac{\Delta\theta}{2}\right) \right]_{-\pi}^\pi \\
    &= \frac{1}{2\pi} \left( 2(1) - 2(-1) \right) = \frac{4}{2\pi} = \frac{2}{\pi} \approx 0.6366.
\end{align}
The expected loss in amplitude is therefore $1 - 2/\pi \approx 36.338\%$. By contrast, on the cylindrical manifold, $A_t = (1-t)A_0 + tA_1$, which trivially evaluates to $1.0$ at $t=0.5$, yielding exactly $0\%$ attenuation.
\end{proof}


\begin{proposition}[Vanishing Gradient of the Chordal Penalty]
\label{prop:vanishing_gradient}
Let the chordal Euclidean loss be $\mathcal{L}_{\mathrm{chord}}(\theta) = \frac{1}{2}\|e^{i\theta} - e^{i\theta^*}\|_2^2 = 1 - \cos(\theta - \theta^*)$ and the geodesic loss be $\mathcal{L}_{\mathrm{geod}}(\theta) = \frac{1}{2}(\theta - \theta^*)^2$. As the phase error reaches its maximum $\Delta\theta = |\theta - \theta^*| \to \pi$, the gradient of the chordal penalty vanishes to zero, while the geodesic gradient reaches its supremum.
\end{proposition}
\begin{proof}
Taking the derivative of the chordal loss with respect to the phase $\theta$:
\begin{equation}
    \frac{\partial \mathcal{L}_{\mathrm{chord}}}{\partial \theta} = \frac{d}{d\theta} (1 - \cos(\theta - \theta^*)) = \sin(\theta - \theta^*).
\end{equation}
Taking the limit as $\Delta\theta \to \pi$, we find $\lim_{\Delta\theta \to \pi} |\sin(\Delta\theta)| = \sin(\pi) = 0$. Thus, at the point of maximum possible error (antipodal phase), the network receives exactly zero corrective gradient signal.
Conversely, the geodesic loss derivative is:
\begin{equation}
    \frac{\partial \mathcal{L}_{\mathrm{geod}}}{\partial \theta} = \theta - \theta^*,
\end{equation}
which linearly increases with the error, providing a maximal restoring force of $\pi$ at $\Delta\theta \to \pi$.
\end{proof}

\begin{proposition}[Angular velocity of the Cartesian bridge]
\label{thm:angular_divergence}
Let $z_0, z_1 \in \mathbb{C}$ with $z_0 \ne z_1$, and let $z_t = (1-t) z_0 + t z_1$. Write $c = \operatorname{Im}(\bar{z}_0 z_1)$. At every $t$ with $z_t \ne 0$ the induced angular velocity is
\begin{equation}
    \dot{\theta}_t = \frac{\operatorname{Im}(\bar{z}_t \dot{z}_t)}{|z_t|^2} = \frac{c}{|z_t|^2},
\end{equation}
so the numerator does not depend on $t$, and
\begin{equation}
    \label{eq:peak_angular}
    \sup_{t \in [0,1]} |\dot{\theta}_t| = \frac{|c|}{d^2}, \qquad d = \min_{t \in [0,1]} |z_t| ,
\end{equation}
with the supremum attained at the closest approach to the origin.
\end{proposition}
\begin{proof}
With $z = x + iy$ and $A = |z|$, the map $\theta = \operatorname{atan2}(y, x)$ is differentiable at every point with $A > 0$, and the chain rule gives $\dot{\theta} = (x\dot{y} - y\dot{x}) / A^2 = \operatorname{Im}(\bar{z}\dot{z}) / |z|^2$. Along the bridge $\dot{z}_t = z_1 - z_0$ is constant, and
\begin{align}
    \operatorname{Im}(\bar{z}_t \dot{z}_t)
    &= \operatorname{Im}\big( ((1-t)\bar{z}_0 + t \bar{z}_1)(z_1 - z_0) \big) \\
    &= (1-t) \operatorname{Im}(\bar{z}_0 z_1) - t \operatorname{Im}(\bar{z}_1 z_0) = \operatorname{Im}(\bar{z}_0 z_1) = c,
\end{align}
because $\operatorname{Im}(\bar{z}_0 z_0) = \operatorname{Im}(\bar{z}_1 z_1) = 0$ and $\operatorname{Im}(\bar{z}_1 z_0) = -c$. The numerator being constant, $|\dot{\theta}_t|$ is largest exactly where $|z_t|$ is smallest.
\end{proof}

Two consequences are worth separating from the usual informal statement. First, the bridge does \emph{not} approach the origin arbitrarily closely for fixed endpoints: combining $|\dot{\theta}| = |c| / A^2$ with the Cauchy--Schwarz bound $|\dot{\theta}| \le |z_1 - z_0| / A$ gives $d \ge |c| / |z_1 - z_0|$, which is the distance from the origin to the line through $z_0$ and $z_1$. Second, an upper bound of order $A^{-1}$ on $|\dot{\theta}|$ does not by itself establish divergence, since a purely radial flow has $c = 0$ and $\dot{\theta} \equiv 0$ while satisfying the same bound. What diverges, and the regime the experiments probe, is the peak over a bridge whose endpoints approach antipodal.

\begin{corollary}[Peak angular velocity at equal amplitudes]
\label{cor:peak_equal_amplitude}
If $|z_0| = |z_1| = A > 0$ and $\delta = \operatorname{wrap}(\theta_1 - \theta_0)$, the closest approach occurs at $t = 1/2$ and
\begin{equation}
    \sup_{t \in [0,1]} |\dot{\theta}_t| = 2 \left| \tan\left( \frac{\delta}{2} \right) \right| ,
\end{equation}
which is unbounded as $\delta \to \pm\pi$ and equals the cylindrical rate $|u_\theta| = |\delta|$ only at $\delta = 0$.
\end{corollary}
\begin{proof}
$|z_t|^2 = A^2\big(1 - 2t(1-t)(1 - \cos\delta)\big)$ is minimised at $t = 1/2$, where $|z_{1/2}| = A|\cos(\delta/2)|$. With $c = A^2 \sin\delta = 2A^2 \sin(\delta/2)\cos(\delta/2)$, equation~\eqref{eq:peak_angular} gives $2|\sin(\delta/2)| / |\cos(\delta/2)|$.
\end{proof}

\begin{corollary}[Heavy tail under a uniform phase difference]
\label{cor:pareto_tail}
If in addition $\delta \sim \mathcal{U}(-\pi, \pi]$, then for $x > 0$
\begin{equation}
    \mathbb{P}\Big( \sup_t |\dot{\theta}_t| > x \Big) = 1 - \frac{2}{\pi} \arctan\left(\frac{x}{2}\right) = \frac{4}{\pi x} + \mathcal{O}(x^{-3}),
\end{equation}
a Pareto tail of index $1$.
\end{corollary}
\begin{proof}
$2|\tan(\delta/2)| > x$ iff $|\delta| > 2\arctan(x/2)$, and $|\delta|$ is uniform on $[0, \pi]$. Expanding $\arctan(x/2) = \pi/2 - 2/x + \mathcal{O}(x^{-3})$ gives the asymptotic form.
\end{proof}

Corollary~\ref{cor:pareto_tail} fixes the amplitudes. When they are random as well, the peak is set by a race between a numerator that vanishes linearly in the smaller amplitude and a squared distance that vanishes quadratically: writing $A_0 = 1$ and $A_1 = a$, both $c$ and $d$ are of order $a |\sin \delta|$, so $\sup_t |\dot{\theta}_t| = |c| / d^2$ is of order $1 / (a |\sin \delta|)$ and small amplitudes enlarge the peak rather than shrink it. The index that results depends on the law of the amplitudes near zero and is not fixed by this argument; an exactly vanishing endpoint contributes nothing at all, since $c = 0$ makes the chord radial and its argument never turns. The index near $1$ reported in Table~\ref{tab:bridge_ablation} is therefore an empirical finding rather than a consequence of this corollary. The cylindrical bridge admits no such tail: $\theta_t = \theta_0 + t u_\theta$ turns at the constant rate $u_\theta$, and $|u_\theta| \le \pi$ for every pair of endpoints.

\section{Both Geometries under Both Losses}
\label{sec:app_losses}
Tables~\ref{tab:unet_baseline} and~\ref{tab:unet_scaling} train both geometries with the same objective, the unweighted squared velocity error. Table~\ref{tab:unet_loss_protocols} repeats the grid with the absolute error as well, so that each geometry can also be compared at its best objective. The best of the four variants is selected on the evaluation itself, so the comparison is post-selection and its nominal $p$-values do not hold; we report it descriptively, as a robustness check rather than as a test.

\begin{table}[t]
    \centering
    \caption{\textbf{Dimensionality Scaling.} Generative error (Sliced $W_2$, $\downarrow$) across spatial resolutions, mean over 5 seeds, with the cylindrical model trained with and without the joint OT coupling. All arms use the unweighted squared velocity error. CyFM's few-step advantage holds at all three resolutions, and joint OT improves the cylindrical model at every step count.}
    \label{tab:unet_scaling}
    \vspace{1mm}
    \resizebox{\linewidth}{!}{
    \begin{tabular}{lllccccc}
        \toprule
        \textbf{Resolution} & \textbf{Geometry} & \textbf{Coupling} & \textbf{1 Step} & \textbf{2 Steps} & \textbf{4 Steps} & \textbf{8 Steps} & \textbf{100 Steps} \\
        \midrule
        $16\times16$ ($D=256$)
        & Cartesian & OT & 0.227 & 0.202 & 0.153 & 0.097 & 0.068 \\
        & CyFM & Independent & 0.116 & 0.155 & 0.133 & 0.097 & 0.062 \\
        & \textbf{CyFM (Ours)} & \textbf{Joint OT} & \textbf{0.098} & \textbf{0.086} & \textbf{0.053} & \textbf{0.064} & 0.054 \\
        \midrule
        $32\times32$ ($D=1024$)
        & Cartesian & OT & 0.340 & 0.275 & 0.209 & 0.104 & 0.037 \\
        & CyFM & Independent & 0.123 & 0.145 & 0.113 & 0.063 & 0.047 \\
        & \textbf{CyFM (Ours)} & \textbf{Joint OT} & \textbf{0.113} & \textbf{0.094} & \textbf{0.072} & \textbf{0.049} & 0.041 \\
        \midrule
        $64\times64$ ($D=4096$)
        & Cartesian & OT & 0.376 & 0.318 & 0.260 & 0.141 & 0.047 \\
        & CyFM & Independent & 0.121 & 0.148 & 0.109 & 0.066 & 0.043 \\
        & \textbf{CyFM (Ours)} & \textbf{Joint OT} & \textbf{0.117} & \textbf{0.114} & \textbf{0.078} & \textbf{0.054} & 0.040 \\
        \bottomrule
    \end{tabular}
    }
    \vspace{1mm}
\end{table}


\begin{table}[h]
    \centering
    \caption{\textbf{Both geometries under both losses.} Generative error (Sliced $W_2$, $\downarrow$), mean over 5 seeds. L1 / L2: absolute / squared error of the velocity; the cylindrical phase term is unweighted in both. Taking each geometry's best of the four variants per cell, CyFM has the lower error at every $k \le 8$ with all seeds separated, except $16\times16$, $k = 8$ ($p = 0.15$); at $k = 100$ no best-against-best comparison is significant ($p \ge 0.54$).}
    \label{tab:unet_loss_protocols}
    \vspace{1mm}
    \resizebox{\linewidth}{!}{
    \begin{tabular}{llllccccc}
        \toprule
        \textbf{Resolution} & \textbf{Geometry} & \textbf{Coupling} & \textbf{Loss} & \textbf{1 Step} & \textbf{2 Steps} & \textbf{4 Steps} & \textbf{8 Steps} & \textbf{100 Steps} \\
        \midrule
        $16\times16$ & Cartesian & Independent & L1 & 0.418 & 0.353 & 0.226 & 0.085 & 0.076 \\
         & Cartesian & Independent & L2 & 0.442 & 0.334 & 0.225 & 0.100 & 0.071 \\
         & Cartesian & OT & L1 & 0.183 & 0.168 & 0.130 & 0.077 & 0.061 \\
         & Cartesian & OT & L2 & 0.227 & 0.202 & 0.153 & 0.097 & 0.068 \\
         & CyFM & Independent & L1 & 0.143 & 0.134 & 0.137 & 0.130 & 0.102 \\
         & CyFM & Independent & L2 & 0.116 & 0.155 & 0.133 & 0.097 & 0.062 \\
         & CyFM & Joint OT & L1 & 0.093 & 0.096 & 0.090 & 0.092 & 0.070 \\
         & CyFM & Joint OT & L2 & 0.098 & 0.086 & 0.053 & 0.064 & 0.054 \\
        \midrule
        $32\times32$ & Cartesian & Independent & L1 & 0.428 & 0.363 & 0.258 & 0.125 & 0.069 \\
         & Cartesian & Independent & L2 & 0.431 & 0.339 & 0.253 & 0.127 & 0.043 \\
         & Cartesian & OT & L1 & 0.320 & 0.263 & 0.199 & 0.091 & 0.048 \\
         & Cartesian & OT & L2 & 0.340 & 0.275 & 0.209 & 0.104 & 0.037 \\
         & CyFM & Independent & L1 & 0.125 & 0.146 & 0.129 & 0.097 & 0.110 \\
         & CyFM & Independent & L2 & 0.123 & 0.145 & 0.113 & 0.063 & 0.047 \\
         & CyFM & Joint OT & L1 & 0.091 & 0.087 & 0.079 & 0.056 & 0.070 \\
         & CyFM & Joint OT & L2 & 0.113 & 0.094 & 0.072 & 0.049 & 0.041 \\
        \midrule
        $64\times64$ & Cartesian & Independent & L1 & 0.411 & 0.365 & 0.283 & 0.136 & 0.037 \\
         & Cartesian & Independent & L2 & 0.416 & 0.349 & 0.285 & 0.154 & 0.046 \\
         & Cartesian & OT & L1 & 0.363 & 0.313 & 0.252 & 0.126 & 0.032 \\
         & Cartesian & OT & L2 & 0.376 & 0.318 & 0.260 & 0.141 & 0.047 \\
         & CyFM & Independent & L1 & 0.107 & 0.169 & 0.145 & 0.091 & 0.050 \\
         & CyFM & Independent & L2 & 0.121 & 0.148 & 0.109 & 0.066 & 0.043 \\
         & CyFM & Joint OT & L1 & 0.092 & 0.122 & 0.099 & 0.064 & 0.035 \\
         & CyFM & Joint OT & L2 & 0.117 & 0.114 & 0.078 & 0.054 & 0.040 \\
        \bottomrule
    \end{tabular}
    }
    \vspace{1mm}
\end{table}


\section{Scaling to Full-Resolution Clinical Matrices: fastMRI Knee at 320\texorpdfstring{$\times$}{x}320}
\label{sec:app_mri_scaling}

We investigate the scaling behaviour of complex generative flows when scaling from $64\times64$ crops to full-resolution clinical MRI matrices ($320\times320$, $D=102{,}400$ complex dimensions). The cohort comprises coronal proton-density (CORPD) knee wavefields from the fastMRI dataset \citep{zbontar2018fastmri}, coil-combined via ESPIRiT into complex fields $z \in \mathbb{C}^{320\times320}$. We train U-Net backbones for 40 epochs across 5 independent seeds under identical training budgets and evaluate generated fields at solver steps $k \in \{1, 2, 4, 8, 100\}$.

Table~\ref{tab:mri_320_scaling} presents the comprehensive 5-seed benchmark comparing CyFM and Cartesian flows across distributional (sliced $W_2$) and spatial texture metrics: lag-1 spatial autocorrelation gap on amplitude ($\text{spatial\_lag1\_gap}$), circular phase lag-1 gap ($\text{phase\_lag1\_gap}$), and the radial power spectral density gap ($\text{radial\_spectrum\_gap}$).

\begin{table}[t]
    \centering
    \caption{\textbf{Scaling to full-resolution clinical matrices ($320\times320$, $D=102{,}400$).} fastMRI Knee CORPD, mean over 5 seeds. Asterisks ($^\star$) indicate separated seeds (5/5). Sliced $W_2$ measures pixelwise distribution fidelity; spatial lag-1, circular phase lag-1, and radial spectrum measure spatial coherence and frequency texture gaps relative to held-out test data ($\downarrow$).}
    \label{tab:mri_320_scaling}
    \vspace{1mm}
    \resizebox{\linewidth}{!}{
    \begin{tabular}{llccccc}
        \toprule
        \textbf{Metric} & \textbf{Arm} & $k=1$ & $k=2$ & $k=4$ & $k=8$ & $k=100$ \\
        \midrule
        \multicolumn{7}{l}{\textit{Sliced $W_2$ ($\downarrow$)}} \\
        & \textbf{Cylindrical + OT (ours)} & \textbf{0.0561}$^\star$ & 0.1228 & 0.1121 & \textbf{0.0821} & 0.0631 \\
        & Cartesian + OT                   & 0.1203 & \textbf{0.1223} & \textbf{0.1166} & 0.0872 & 0.0720 \\
        & Cartesian, independent           & 0.1349 & 0.1353 & 0.1256 & 0.0872 & \textbf{0.0629} \\
        \midrule
        \multicolumn{7}{l}{\textit{Spatial Lag-1 Gap ($\downarrow$)}} \\
        & \textbf{Cylindrical + OT (ours)} & 0.5808 & 0.3444 & 0.1295 & 0.0616 & 0.0321 \\
        & Cartesian + OT                   & \textbf{0.1532}$^\star$ & \textbf{0.0051}$^\star$ & \textbf{0.0072}$^\star$ & \textbf{0.0083}$^\star$ & \textbf{0.0014}$^\star$ \\
        \midrule
        \multicolumn{7}{l}{\textit{Circular Phase Lag-1 Gap ($\downarrow$)}} \\
        & \textbf{Cylindrical + OT (ours)} & 0.7366 & 0.3331 & 0.2304 & 0.1655 & 0.1190 \\
        & Cartesian, independent           & \textbf{0.0941}$^\star$ & \textbf{0.0297}$^\star$ & \textbf{0.0521}$^\star$ & \textbf{0.0733}$^\star$ & \textbf{0.0671}$^\star$ \\
        \midrule
        \multicolumn{7}{l}{\textit{Radial Spectrum Gap ($\downarrow$)}} \\
        & \textbf{Cylindrical + OT (ours)} & 1.5877 & 1.2927 & 0.8364 & 0.5732 & 0.3882 \\
        & Cartesian, independent           & \textbf{1.2978}$^\star$ & \textbf{0.2944}$^\star$ & \textbf{0.1372}$^\star$ & \textbf{0.1662}$^\star$ & \textbf{0.0706}$^\star$ \\
        \bottomrule
    \end{tabular}
    }
    \vspace{1mm}
\end{table}

\textbf{What the benchmark shows, and what it does not.}
\begin{enumerate}
    \item \textbf{The single-step advantage survives the dimension.} At $k=1$ CyFM attains a
    sliced $W_2$ of $0.0561$, $2.1\times$ lower than Cartesian$+$OT ($0.1203$) and $2.4\times$
    lower than Cartesian independent ($0.1349$), with all five seeds separated. Bounding the
    angular target by $\pi$ therefore still buys a single-step advantage at $102{,}400$
    dimensions.
    \item \textbf{Every spatial measure prefers the plane, at every step count.} This is the
    reverse of the $64\times64$ result, and we report it as measured. The Cartesian arm's
    lag-one coherence does not fall short of the reference but exceeds it, so its advantage on
    these gaps is the advantage of a smoother field than the data rather than a more faithful
    one; the gap statistic already charges it for overshooting, and it is still closer.
    \item \textbf{The coupling is not the cause.} A cylindrical arm trained with no coupling at
    all reaches $0.3926$ on the radial spectrum gap against $0.3882$ for the OT-coupled arm, and
    a Cartesian arm trained \emph{with} minibatch OT at the same $102{,}400$ dimensions reaches
    $0.0756$. The collapse of the transport-cost reduction with dimension
    (Section~\ref{sec:exp_scaling}) therefore does not explain this, and cannot be offered as an
    explanation: the failure is a property of the coordinate representation, not of the pairing.
    \item \textbf{The one measured difference between the two matrices.} At $320\times320$,
    $16.7\%$ of coefficients are exactly zero, and $\operatorname{atan2}(0,0)$ assigns each of
    them a phase of $0$ --- a placeholder rather than a measurement. The $64\times64$ crop
    contains no exactly zero coefficient. Under the decoupled metric the phase term is weighted
    equally everywhere, so at $320\times320$ roughly one pixel in six contributes a regression
    target that carries no information, while under the inherited metric the factor $A^2$
    suppresses exactly those pixels. Consistent with this, weighting the phase term by the
    amplitude improves the cylindrical phase gap at $320\times320$ ($0.119 \to 0.068$) and
    \emph{worsens} it on the $64\times64$ crop ($0.106 \to 0.166$), where there is no such
    background to suppress. We offer this as the candidate the evidence supports, not as a
    demonstration: it predicts the sign of the weighting effect at both matrices, but it has not
    been shown to account for the size of the spatial gaps.
\end{enumerate}
Separating a genuine non-measurement from a small measurement is therefore the concrete obstacle
to scaling this geometry to uncropped clinical matrices, and a mask derived from the acquisition
rather than from an amplitude threshold is the form a fix would take. We leave it open.

\section{The Speech Grid at Two Training Budgets}
\label{sec:app_audio_epochs}
The speech block of Table~\ref{tab:unet_baseline} was first run for two epochs, and at that
budget the Cartesian arm had the lower error at $k=100$. The schedule explains it rather than
the geometry: the cosine annealing is parameterised by the epoch count, so a two-epoch run
drives the learning rate to its floor within those two epochs and trains almost nothing. The
training loss was still falling when it stopped. Table~\ref{tab:audio_epochs} repeats the grid
at ten and at forty epochs, where the schedule completes; neither reproduces the two-epoch
ordering, and the cylindrical arms move little between the two budgets while the Cartesian
arms do not overtake them at either.

\begin{table}[h]
    \centering
    \caption{\textbf{The speech block is not an artefact of the training budget.} The same
    grid at 10 and at 40 epochs, mean over 5 seeds. The cylindrical arms move little with
    the budget; the Cartesian arms do not overtake them at either. A two-epoch probe, run
    first, put the Cartesian arm ahead at $k=100$ ($0.0210$): its cosine schedule had
    annealed the learning rate to its floor within those two epochs, so that model was
    barely trained, and neither longer run reproduces it.}
    \label{tab:audio_epochs}
    \vspace{1mm}
    \begin{tabular}{llccccc}
        \toprule
        \textbf{Epochs} & \textbf{Arm} & $k=1$ & $k=2$ & $k=4$ & $k=8$ & $k=100$ \\
        \midrule
        10
          & Cylindrical + OT       & 0.0461 & 0.0435 & 0.0371 & 0.0348 & 0.0292 \\
          & Cylindrical, indep.    & 0.0467 & 0.0440 & 0.0377 & 0.0348 & 0.0272 \\
          & Cartesian + OT         & 0.0766 & 0.0687 & 0.0563 & 0.0471 & 0.0407 \\
          & Cartesian, indep.      & 0.0803 & 0.0748 & 0.0608 & 0.0512 & 0.0455 \\
        \midrule
        40
          & Cylindrical + OT       & 0.0454 & 0.0403 & 0.0333 & 0.0308 & 0.0264 \\
          & Cylindrical, indep.    & 0.0458 & 0.0409 & 0.0336 & 0.0303 & 0.0261 \\
          & Cartesian + OT         & 0.0780 & 0.0695 & 0.0563 & 0.0442 & 0.0353 \\
          & Cartesian, indep.      & 0.0822 & 0.0758 & 0.0610 & 0.0457 & 0.0353 \\
        \bottomrule
    \end{tabular}
    \vspace{1mm}
\end{table}
